\documentclass[11pt,a4paper]{article}
\usepackage{times}
\usepackage{latexsym}
\usepackage{xcolor}
\usepackage{url}
\usepackage{booktabs}
\usepackage{amsmath}

\newcommand\BibTeX{B{\sc ib}\TeX}

\begin{document}

\title{Persona-Aware Voice Assistant for Structured Tool Decisions}

\date{}
\maketitle

\begin{abstract}
This project builds a persona-aware listening assistant that maps natural spoken dialogue into structured tool decisions. Given an utterance, the model predicts an importance label, a tool call, and a short detail string as JSON. The system supports creating reminders, calendar events, alarms, and notes, as well as modifying or deleting existing items, asking for clarification when needed, and deciding when no action should be taken. I focus on four ideas: adding an explicit system prompt that explains how to decide, constraining the output to valid JSON using three complementary methods, comparing current-utterance-only prediction with previous-two-turn context, and adding explicit saved-state context for edit/delete decisions. I evaluate both an Arjun-only model and an all-persona model trained across Arjun, Dev, Margaret, and Neel while conditioning each example on the correct persona prompt. On the Arjun held-out split, the state-plus-context model reaches 0.492 tool accuracy, 0.574 importance accuracy, 0.377 exact-match accuracy, 1.000 JSON parse rate, and 0.885 action F1. On the all-persona held-out split, the best exact-match result comes from the state-plus-context model, with 0.637 tool accuracy, 0.689 importance accuracy, 0.479 exact-match accuracy, 0.984 JSON parse rate, and 0.928 action F1.
\end{abstract}

\section{Project Overview (10 points)}

Voice assistants are increasingly useful for reminders, notes, scheduling, and personal task management. However, a useful assistant must do more than transcribe speech: it must decide whether an utterance requires an action, which tool should be called, and how urgent or important that action is. This is difficult because conversational language is indirect. A user might say ``don't let me forget,'' ``put that on my calendar,'' ``remember that,'' or ``actually cancel it,'' and each phrase implies a different tool decision. The goal of this project is to build a persona-aware assistant that converts dialogue into a structured JSON decision:
\begin{verbatim}
{"importance": "Low|Medium|High",
 "tool": "tool name or null",
 "detail": "string"}
\end{verbatim}
The input is a spoken or written utterance, optionally with prior dialogue context. The output is a machine-readable decision that can be used by a downstream tool executor.

Existing language models are strong at generating plausible text, but they are not automatically reliable as tool routers. They may confuse a reminder with a note, a calendar event with a reminder, or an alarm with a generic future task. They may also produce invalid JSON, extra explanations, or unsupported tool names. Existing assistants often rely on generic intent recognition or prompt-only instruction following, but this is not enough when decisions must be personalized. For example, the meaning of a calendar event is universal, but the importance of a sports practice, medication refill, or research meeting depends on the user. Therefore, the system needs both stable tool semantics and persona-specific importance prediction.

I propose four ideas to address this gap. First, I add a system prompt that explains the decision process: tool meanings are universal and persona-independent, while importance is persona-dependent. Second, I constrain the output format with three layers: strict JSON prompting, post-generation parsing and validation, and guided structured decoding for the live assistant. Third, I test whether adding the previous two dialogue turns improves decision quality compared with using only the current utterance. Fourth, I add an explicit saved-state representation so that the model can see existing reminders, notes, alarms, and calendar events before deciding whether an utterance should create, update, delete, clarify, or do nothing. I run these ideas in both a single-persona Arjun setting and an all-persona setting that combines Arjun, Dev, Margaret, and Neel.

Although dataset improvement is not treated as a separate idea section, it was still an important implementation step. Weak classes such as Medium importance, alarms, updates, deletes, no-action examples, and clarification requests were difficult for the model. I added targeted persona-matched examples that distinguish ``reserve time'' from ``do not forget,'' urgent alarms from ordinary reminders, reference notes from future tasks, and ambiguous requests from executable tool calls.

The system prompt addresses ambiguity in the task definition. It gives the model a consistent decision order: first choose the tool from universal definitions, then choose importance using the persona's goals, risks, habits, and deadlines. This matters because tool semantics should not change across people, but importance should.

The JSON constraints address deployment reliability. A model can make a reasonable semantic decision but still fail if the output is not parseable. I therefore implemented prompt-level format instructions, a validator that checks the generated JSON schema, and structured decoding through Ollama for the live voice path.

I evaluate the ideas using held-out dialogue examples. The main questions are: Does supervised fine-tuning improve tool and importance decisions over the base model? Does the model produce valid JSON? Which weak classes remain difficult? Does previous dialogue context improve decision quality? I measure tool accuracy, importance accuracy, action precision, action recall, action F1, JSON parse rate, exact-match accuracy, per-tool precision/recall/F1, confusion matrices, and weak-class recall.

The main finding is that all-persona training improves overall robustness, while saved state is especially useful for state-dependent actions such as updates and deletes. In the Arjun-only setting, state-plus-context is clearly strongest. In the all-persona setting, the current-only SFT model has slightly higher tool accuracy, but state-plus-context has the best importance accuracy, exact-match accuracy, and action recall.

\section{Ideas (10 points)}

\subsection{Idea 1: Adding a System Prompt on How to Decide}

The first idea was to add a system prompt that explains the model's decision process. The prompt states that tool semantics are universal, but importance is persona-specific. This is the central architecture of the project.

Universal tool definitions are:
\begin{itemize}
    \item Reminder: use when the user should not forget a future task.
    \item Calendar event: use when time should be reserved.
    \item Alarm: use for an urgent prompt at a specific moment.
    \item Note: use for reference material worth saving.
    \item Update/delete: use when an existing item should be changed or removed.
    \item Null: use when no action is needed.
\end{itemize}

Importance, however, depends on the persona. A medication refill may be high importance for Margaret, while a research paper idea or manuscript task may be high importance for Neel because he is an enthusiastic neuroscience researcher and these tasks directly support his academic goals. A looser paper-related errand, such as sending an article to a family member, may be medium instead. A tournament registration may be high for Arjun if it has a deadline, while a casual sports joke should be low. The prompt tells the model to choose the tool first using the universal rules, then choose importance using the user's goals, habits, deadlines, and risks.

This idea is applicable because the model needs a stable policy, not only examples. Its strength is interpretability: the decision rules can be inspected and reused across datasets. The challenge is that the prompt cannot guarantee the model follows the rules.

\subsection{Idea 2: Constraining JSON Format}

The second idea was to make the output reliably machine-readable. I implemented three JSON-format constraints.

\paragraph{Approach 1: strict JSON prompt.}
The system prompt and dataset targets instruct the model to return exactly one JSON object with the keys \texttt{importance}, \texttt{tool}, and \texttt{detail}. It also forbids markdown, explanations, comments, and code fences. This teaches the model the desired format during supervised training.

\paragraph{Approach 2: parser and validator.}
I implemented the post-generation parser and validator in \texttt{json\_decision.py}. After the model generates text, this utility first extracts the first complete JSON object from the raw output. It then checks whether the object follows the required decision format: \texttt{importance} must be one of \texttt{Low}, \texttt{Medium}, or \texttt{High}; \texttt{tool} must be either \texttt{null} or one of the allowed tool names; and \texttt{detail} must be a string. The validator also normalizes small formatting issues, such as treating an empty tool value as \texttt{null} and using an empty detail string when the tool is \texttt{null}. If the output cannot be parsed or does not satisfy these constraints, the system treats it as invalid and can either attempt repair or fall back to a safe null decision. This approach validates model outputs programmatically; it is not a separate validation dataset.

\paragraph{Approach 3: guided structured decoding.}
For the live Ollama voice path, I pass a JSON schema to the model using structured decoding. This constrains the generation process toward valid JSON, rather than only checking after generation. The schema path is the strongest option because it specifies the exact fields and allowed values. If schema mode is unavailable, the code first tries a weaker JSON mode that asks the model to produce valid JSON, and if that is also unavailable, it uses normal generation and then applies the parser and validator from \texttt{json\_decision.py}.

These three approaches solve different parts of the problem. Prompting teaches the behavior, validation detects failures, and guided decoding prevents many format failures during generation. A key observation is that strict prompting plus validation can still produce JSON parse rate below 1.0 because Hugging Face evaluation uses free-form generation. Guided decoding is stronger, but it is currently used in the live voice/Ollama path rather than the Hugging Face evaluation path.

\subsection{Idea 3: Without Context vs. With Previous Two Turns}

The third idea was to compare current-utterance-only prediction against prediction with previous dialogue context. Many utterances are underspecified without context. For example, ``put that on my calendar'' requires knowing what ``that'' refers to. Similarly, ``actually move it to tomorrow'' or ``delete that reminder'' needs previous dialogue.

I implemented a context-window dataset where each example includes the previous two dialogue turns plus the current utterance:
\begin{verbatim}
Recent conversation:
Speaker 1: ...
Speaker 2: ...

Current utterance:
Speaker 1: ...
\end{verbatim}


This tests whether the model improves when exposed to the recent dialogue context.

\subsection{Idea 4: State Plus Context}

The fourth idea was to add explicit saved-state context. Previous turns help resolve pronouns, but edit/delete decisions often require more than recent dialogue. If a user says ``move that reminder to tomorrow morning,'' the assistant should know what reminder currently exists. If the user says ``delete the alarm,'' the assistant should know whether an alarm is already saved.

I therefore created a state-plus-context dataset. Each input may include:
\begin{verbatim}
Recent conversation:
Speaker 1: ...
Speaker 2: ...

Existing saved items:
- Reminder: ...
- Calendar: ...
- Alarm: ...
- Note: ...

Current utterance:
Speaker 1: ...
\end{verbatim}

The saved-state list is built from earlier tool calls in the same conversation. Create calls add state, update calls modify state, delete calls remove state, and null calls leave state unchanged. This representation is closer to an actual tool-using assistant, because the model is no longer deciding from language alone; it can condition on the user's memory.


\section{Experimental Setup (10 points)}

\subsection{Models}

The base model is Qwen2.5-1.5B-Instruct, a pretrained transformer instruction-following model. I fine-tuned it using supervised fine-tuning with LoRA through Unsloth. I used Qwen2.5-1.5B because it is small enough to train on a Colab T4 GPU while still being strong enough for instruction following.

The main training settings were:
\begin{itemize}
    \item Model: Qwen2.5-1.5B-Instruct.
    \item Fine-tuning method: supervised fine-tuning with LoRA.
    \item LoRA rank: 32.
    \item Learning rate: $2 \times 10^{-4}$.
    \item Epochs: 10.
    \item Batch size per device: 2.
    \item Gradient accumulation steps: 4.
    \item Precision: bf16 when supported, otherwise fp16 on T4.
\end{itemize}

I trained three supervised variants for the main experiment: a current-only SFT model, a previous-two-turn context SFT model, and a state-plus-previous-two-turn context SFT model.

\subsection{Dataset}

The dataset consists of persona-specific conversations for Arjun, Dev, Margaret, and Neel. Each persona has 18 conversations. Conversations 9, 10, 16, and 17 are held out for testing, and the remaining 14 conversations are used for training. Conversation 18 stays in training because it contains targeted weak-class examples such as update, delete, alarm, no-action, and clarification cases.

For Arjun:
\begin{itemize}
    \item Training conversations: 14.
    \item Test conversations: 4.
    \item Training rows: 254.
    \item Test rows: 61.
\end{itemize}

For the all-persona experiment:
\begin{itemize}
    \item Training rows: 779.
    \item Test rows: 190.
    \item Training rows by persona: 254 Arjun, 222 Dev, 157 Margaret, and 146 Neel.
    \item Test rows by persona: 61 Arjun, 55 Dev, 39 Margaret, and 35 Neel.
\end{itemize}

The label format is the same for both current-only and context training:
\begin{verbatim}
{"importance": "Low|Medium|High",
 "tool": "allowed tool name or null",
 "detail": "string"}
\end{verbatim}
The difference is the input. The current-only dataset contains only the current utterance. The context dataset includes the previous two turns plus the current utterance. The state-plus-context dataset includes the previous two turns, a list of existing saved items, and the current utterance.

\subsection{Evaluation Metrics}

I use automatic evaluation against the gold labels in the held-out conversations.

\paragraph{Tool accuracy.}
The fraction of examples where the predicted tool equals the gold tool.

\paragraph{Importance accuracy.}
The fraction of examples where the predicted importance equals the gold importance.

\paragraph{Action precision, recall, and F1.}
For this metric, any non-null tool is treated as an action. Precision is the fraction of predicted actions that correspond to gold actions. Recall is the fraction of gold actions detected by the model. F1 is:
\[
F1 = \frac{2PR}{P+R}
\]

\paragraph{JSON parse rate.}
The fraction of outputs that can be parsed and validated as schema-compatible JSON.

\paragraph{Exact-match accuracy.}
The fraction of examples where both tool and importance are correct.

\paragraph{Detail quality.}
I also evaluate the generated \texttt{detail} argument using exact string match, whether the model correctly predicts the presence or absence of a detail, and token-level F1. Exact string match is strict, while token-level F1 better captures partial semantic overlap.

\paragraph{Diagnostic metrics.}
I also compute tool confusion matrices, importance confusion matrices, action-vs-no-action confusion, per-tool precision/recall/F1, and weak-class recall. These diagnostics help decide what data to add next.

\section{Results (40 points for implementation and results)}

Table~\ref{tab:arjun-results} shows the latest Arjun experiment on the four held-out conversations. The strongest method is the model trained and evaluated with previous-two-turn context plus explicit saved-state context.

\begin{table}[t!]
\centering
\small
\begin{tabular}{lrrrrrrrr}
\toprule
Method & Parse & Tool & Imp. & Exact & Prec. & Rec. & F1 & Detail F1 \\
\midrule
Base, current test & 1.000 & 0.279 & 0.475 & 0.148 & 0.804 & 0.957 & 0.874 & 0.228 \\
SFT, current test & 0.262 & 0.262 & 0.262 & 0.213 & 0.800 & 0.170 & 0.281 & 0.231 \\
Ctx2 SFT, context test & 0.459 & 0.311 & 0.311 & 0.197 & 0.821 & 0.489 & 0.613 & 0.182 \\
State+Ctx2 SFT, context test & \textbf{1.000} & \textbf{0.492} & \textbf{0.574} & \textbf{0.377} & \textbf{0.807} & \textbf{0.979} & \textbf{0.885} & \textbf{0.392} \\
\bottomrule
\end{tabular}
\caption{Arjun held-out test results. Ctx2 means previous two dialogue turns plus the current utterance. State+Ctx2 additionally includes existing saved reminders, calendar events, alarms, and notes.}
\label{tab:arjun-results}
\end{table}

The state-plus-context model improves substantially over the current-only SFT and the previous-two-turn context model. Compared with the Ctx2 SFT model, tool accuracy increases from 0.311 to 0.492, importance accuracy increases from 0.311 to 0.574, exact-match accuracy increases from 0.197 to 0.377, JSON parse rate increases from 0.459 to 1.000, action F1 increases from 0.613 to 0.885, and detail token F1 increases from 0.182 to 0.392.

Table~\ref{tab:all-results} shows the all-persona experiment. This model is trained on Arjun, Dev, Margaret, and Neel together, while each example still uses the correct persona-specific system prompt. The all-persona SFT models are much stronger than the base model. The current-only SFT model has the highest tool accuracy at 0.653, but the state-plus-context model has the best importance accuracy, exact-match accuracy, and action recall. The Ctx2 model has the highest detail token F1.

\begin{table}[t!]
\centering
\small
\begin{tabular}{lrrrrrrrr}
\toprule
Method & Parse & Tool & Imp. & Exact & Prec. & Rec. & F1 & Detail F1 \\
\midrule
Base, current test & 0.989 & 0.274 & 0.474 & 0.163 & 0.770 & 0.944 & 0.848 & 0.280 \\
SFT, current test & 0.995 & \textbf{0.653} & 0.589 & 0.447 & \textbf{0.921} & 0.908 & 0.915 & 0.482 \\
Ctx2 SFT, context test & \textbf{1.000} & 0.637 & 0.653 & 0.447 & 0.906 & 0.951 & \textbf{0.928} & \textbf{0.538} \\
State+Ctx2 SFT, context test & 0.984 & 0.637 & \textbf{0.689} & \textbf{0.479} & 0.901 & \textbf{0.958} & \textbf{0.928} & 0.500 \\
\bottomrule
\end{tabular}
\caption{All-persona held-out test results across Arjun, Dev, Margaret, and Neel. The combined test set has 190 examples.}
\label{tab:all-results}
\end{table}

The all-persona results show that adding more persona-matched training data helps substantially. Compared with the base model, the current-only all-persona SFT model improves tool accuracy from 0.274 to 0.653 and exact-match accuracy from 0.163 to 0.447. Adding context and state does not further improve tool accuracy in this run, but it does improve importance accuracy and exact match. The State+Ctx2 model reaches 0.689 importance accuracy and 0.479 exact-match accuracy, which are the best overall decision-level results in the all-persona experiment.

The diagnostic metrics show why state matters. On Arjun, the State+Ctx2 model has an action matrix of 46 true positives, 11 false positives, 1 false negative, and 3 true negatives. This means it detects nearly all actionable utterances while still preserving some no-action behavior. It also improves several update/delete classes: \texttt{update\_calendar\_event} reaches F1 0.889, \texttt{update\_alarm} reaches F1 0.750, and \texttt{delete\_reminder} reaches F1 0.667. These are exactly the kinds of decisions that benefit from knowing what items already exist.

On the all-persona test set, State+Ctx2 has an action matrix of 136 true positives, 15 false positives, 6 false negatives, and 33 true negatives. It performs well on several weak classes: \texttt{update\_calendar\_event} reaches F1 0.625, \texttt{update\_alarm} reaches F1 0.889, \texttt{delete\_reminder} reaches F1 0.815, and \texttt{create\_alarm} reaches F1 0.571. The combined dataset therefore helps rare tool classes, even though clarification remains difficult.

One remaining weakness is clarification. The Arjun test set contains one \texttt{ask\_clarification} example, and none of the evaluated models predicted it correctly. The all-persona test set contains four clarification examples; the current-only SFT model predicts one correctly, while the context and state-plus-context models still miss this class. This suggests that clarification needs more direct training support and possibly a higher-level policy rule: when the user asks to edit or delete an underspecified item, the assistant should ask for clarification rather than guessing.

\section{Analysis and Discussion (25 points)}

The first major failure mode is confusion between similar tools. The base model often detects that an utterance is actionable, but it does not reliably choose the exact tool. This is why its action F1 is high while its tool accuracy is low. For example, on the all-persona test set, the base model has 0.848 action F1 but only 0.274 tool accuracy. In practice, this is not enough for a tool-using assistant: calling a reminder when the user wanted a calendar event still creates the wrong external action.

The second failure mode is Medium importance. Medium is a boundary class: it means the item is useful and actionable, but not urgent enough to be High. In the Arjun result, Ctx2 SFT has Medium recall 0.545 and State+Ctx2 has Medium recall 0.455. This is better than the base model, which has 0.000 Medium recall, but it remains a difficult class. The all-persona model improves importance accuracy overall, with State+Ctx2 reaching 0.689, suggesting that broader persona coverage helps the model learn the boundary between Low, Medium, and High.

The third failure mode is no-action versus action. A safe assistant should not create tools for vague or canceled plans, but it also should not ignore real tasks. The Arjun State+Ctx2 model has very high action recall, with only one false negative, but it still has eleven false positives. The all-persona State+Ctx2 model shows the same pattern at larger scale, with 136 true positives, 15 false positives, 6 false negatives, and 33 true negatives. This means the model is eager to act. In a real assistant, this should be balanced by stronger clarification behavior before risky updates or deletes.

The JSON-format work addresses a different class of failures. Even if the predicted decision is semantically close, the assistant cannot execute it if the output is malformed. Prompt-only constraints improve behavior but do not guarantee parseable JSON. Parser validation makes evaluation more reliable, and guided decoding is the best option for deployment because it constrains generation itself.

Overall, the results suggest that the best architecture is: persona-matched weak-class data, explicit universal tool definitions, persona-aware importance instructions, strict JSON validation, and matched state/context training and inference. The all-persona experiment also shows that combining personas can improve rare-class coverage and general robustness while still preserving persona-specific importance through the system prompt. The most important lesson is that edit and delete tool use is not just a text classification problem. It is a state-conditioned decision problem.

\section{Code and Demo (5 points)}

The main files are:
\begin{itemize}
    \item \texttt{parse\_json\_dataset.py}: creates current-only, context-window, and state-plus-context JSONL datasets.
    \item \texttt{train\_sft.py}: trains the supervised fine-tuned model.
    \item \texttt{train\_sft\_ctx2.py}: trains the previous-two-turn context SFT model.
    \item \texttt{train\_sft\_state\_ctx2.py}: trains the state-plus-context SFT model.
    \item \texttt{build\_combined\_persona\_dataset.py}: builds all-persona training and test files.
    \item \texttt{eval.py}: evaluates base, current SFT, context SFT, and state-plus-context SFT models.
    \item \texttt{json\_decision.py}: implements JSON schema, parser, validator, repair, and structured decoding helpers.
    \item \texttt{colab\_retrain\_one\_persona.ipynb}: trains and evaluates one persona.
    \item \texttt{colab\_retrain\_all\_personas.ipynb}: trains and evaluates all personas.
    \item \texttt{voice\_ui.py}: safe voice demo that logs model decisions.
    \item \texttt{listening\_agent.py}: live assistant path that can call local tools.
\end{itemize}

Example commands:
\begin{verbatim}
python3 parse_json_dataset.py arjun_...json arjun --holdout 9 10 16 17
python3 parse_json_dataset.py arjun_...json arjun --holdout 9 10 16 17 \
  --context-turns 2 --suffix ctx2
python3 parse_json_dataset.py arjun_...json arjun --holdout 9 10 16 17 \
  --context-turns 2 --include-state --suffix state_ctx2

PERSONA=arjun python3 train_sft.py
PERSONA=arjun python3 train_sft_ctx2.py
PERSONA=arjun python3 train_sft_state_ctx2.py

PERSONA=arjun python3 eval.py \
  --base Qwen/Qwen2.5-1.5B-Instruct \
  --sft ./arjun-sft-merged \
  --ctx-sft ./arjun-ctx2-sft-merged \
  --state-ctx-sft ./arjun-state-ctx2-sft-merged \
  --test ./arjun_state_ctx2_test.jsonl
\end{verbatim}

The main software requirements are Python, PyTorch, Transformers, TRL, Unsloth, datasets, peft, bitsandbytes, Ollama, Flask, faster-whisper, and sounddevice. The Colab notebooks contain the full training and evaluation workflow.

\section{Learning Outcomes}

I learned that structured assistant behavior depends heavily on both data quality and input representation. The most valuable examples are not only straightforward commands, but confusing boundary cases that teach the difference between similar tools. I also learned that JSON validity and semantic correctness are separate problems. A model can understand the intent but still produce unusable output if the JSON is invalid. Finally, I learned that dialogue context alone is not enough for edit/delete behavior; explicit saved state makes the decision problem much closer to the real assistant setting.

\end{document}
